\section{Formal checks on rational-thinking improvement examples}
\label{sec:p19-annex-rationality-15-3}

\begin{annexitem}{annex:p19.r153.example-checks}{thm:p19-r153-example-checks}{Validity, relevant cards, irrecoverable cost and screening}
\begin{proof}
For the weather argument take \(P\) false and \(Q\) true.
Then \(P\to Q\) and \(\neg P\) are true, while
\(\neg Q\) is false. This supplies a countermodel to the claimed
entailment under the declared material-conditional semantics. Wet streets
from another source give the corresponding ordinary interpretation.

For the cards, a violation has a vowel on its letter side and an odd
number on its number side. Turning A over can disclose an odd number;
turning 7 over can disclose a vowel. After inspecting these two, neither
remaining card can conceal such a violation: B is already a nonvowel,
and 4 already an even number. Hence these two inspections suffice.
For any fixed inspection set omitting A, choose the hidden letter of 7
to be a consonant and complete all other cards so they satisfy the rule.
The unseen number of A can then be even or odd while the inspected
observations stay fixed, making the entire four-card rule true or false.
For a fixed set omitting 7, make A's hidden number even and complete
the other cards consistently; 7's hidden letter can be a consonant or
a vowel with the same observed inspections, again changing the rule's
status. Hence every fixed set sufficient for all completions must contain
both A and 7. This ex ante requirement does not force further checking
after a violation has already been observed: such a violation immediately
determines that the rule fails.

Subtracting the same real number \(K\) preserves a strict ordering:
the difference between \(u_W-K\) and \(u_M-K\) is
\(u_W-u_M>0\). The past expenditure is the same under either
choice and cannot be recovered. For a complete counterexample to an
unqualified prescription, take \(u_W=10,u_M=9\) before an additional
Wisconsin-specific disutility of 2. Its total becomes 8, below 9.
This changes the future utility premise; it does not make \(K\)
recoverable or justify comparing the face prices as future returns.

For screening, the two disjoint positive-test cases have probabilities
\(.01(.90)=.009\) and \(.99(.09)=.0891\). Their sum
\(.0981\) is positive, so conditional probability is defined and
\[
 \mathbb P(C\mid +)=\frac{.009}{.0981}
   =\frac{90}{981}=\frac{10}{109}.
\]
The displayed frequency construction realizes the same proportions:
10000 times the prevalence gives 100 cancer cases, ninety percent
of which are positive; the remaining 9900 people yield nine percent,
or 891, positives. There are consequently 981 positives, 90 with
cancer. Dividing gives approximately .0917431. Rounding this to one
decimal probability yields .1, but treating .1 as the exact answer
would alter the stipulated numerical calculation. The construction
assumes the given rates, not additional clinical observations.
\end{proof}
\end{annexitem}
